\chapter{Aksjomat wyboru i jego równoważne postacie}\label{ch:axiom-choice}
\chaptermark{Aksjomat wyboru i jego równoważne postacie}

\section{Wybrać jeden element czy wybierać nieskończenie wiele razy?}
Z jednego niepustego zbioru $A$ możemy wskazać element $a\in A$;
nie wymaga to osobnego aksjomatu. Co innego, jeśli mamy rodzinę
$(A_i)_{i\in I}$ niepustych zbiorów i chcemy \emph{jednej funkcji},
która dla każdego $i\in I$ podaje element $A_i$.
Nawet jeśli każde żądanie z osobna da się spełnić,
nie znamy jeszcze reguły gwarantującej dokonanie
wszystkich wyborów jednocześnie.

\begin{definicja}[Aksjomat wyboru, AC]\label{def:choice}
Dla każdego zbioru indeksów $I$ i każdej rodziny
$(A_i)_{i\in I}$ zbiorów niepustych istnieje funkcja
$c:I\to\bigcup_{i\in I}A_i$ taka, że
\[
 \forall i\in I\quad c(i)\in A_i.
\]
Funkcję $c$ nazywamy \emph{funkcją wyboru} tej rodziny.
Teoria $\mathrm{ZFC}$ to $\mathrm{ZF}+\mathrm{AC}$.
\end{definicja}

\begin{przyklad}[Dlaczego rodzina skończona nie wymaga AC?]
Dla dwóch niepustych zbiorów $A_0,A_1$ istnieją elementy
$a_0\in A_0$ i $a_1\in A_1$, więc funkcja
$c(0)=a_0$, $c(1)=a_1$ jest funkcją wyboru.
Przez zwykłą indukcję po liczbie zbiorów otrzymamy tę samą
własność dla każdej rodziny \emph{skończonej}.
AC jest nowym postulatem dopiero w przypadku rodzin dowolnych,
niekoniecznie skończonych.
\end{przyklad}

\begin{przyklad}[Kiedy wybór jest dany konkretnym wzorem?]
Jeżeli wszystkie $A_i$ są niepustymi podzbiorami $\N$,
możemy wybrać $c(i)=\min A_i$. Minimum istnieje dzięki
zwykłemu dobremu uporządkowaniu $\N$.
Taką rodzinę obsługujemy bez dodatkowego aksjomatu.
Nie oznacza to jednak, że dowolny zbiór ma z góry
wyróżniony najmniejszy element: nie każdy zbiór ma
wcześniej ustalony dobry porządek.
\end{przyklad}

\section{Dobre uporządkowanie i jego początkowe odcinki}
\begin{definicja}[Twierdzenie o dobrym uporządkowaniu]
\label{def:wellorder-principle}
Przez \emph{zasadę dobrego uporządkowania} rozumiemy zdanie:
\emph{każdy zbiór $A$ można wyposażyć w porządek liniowy $\le_A$,
w którym każdy niepusty podzbiór ma element najmniejszy}.
Dla dobrego porządku $<$ na $W$ zbiór
$W_{<w}:=\{v\in W:v<w\}$ nazywamy \emph{właściwym odcinkiem
początkowym} określonym przez $w\in W$.
\end{definicja}

\begin{lemat}[Dobre porządki można jednoznacznie porównać]
\label{lem:choice-wellorder-comparison}
Niech $(U,<_{U})$ i $(V,<_{V})$ będą zbiorami dobrze
uporządkowanymi. Wówczas dokładnie jedna z trzech
możliwości zachodzi: są izomorficzne jako zbiory uporządkowane;
$U$ jest izomorficzny z właściwym odcinkiem początkowym $V$;
$V$ jest izomorficzny z właściwym odcinkiem początkowym $U$.
\end{lemat}
\begin{proof}
\emph{Odcinkiem początkowym} nazywamy także cały zbiór;
dla odcinka właściwego wymagamy, aby był postaci $U_{<u}$.
Najpierw zauważmy, że każdy izomorfizm porządkowy między
odcinkami początkowymi jest jednoznaczny: dwa takie
izomorfizmy nie mogą różnić się na wspólnej dziedzinie.
Gdyby się różniły, wybierzmy najmniejszy argument $x$,
na którym dają różne wyniki. Oba przekształcenia muszą
wtedy przenosić zbiór wcześniejszych elementów $U_{<x}$
na zbiór wcześniejszych elementów od odpowiednich obrazów.
Ponieważ na $U_{<x}$ są takie same, mają ten sam obraz
odcinka, a w dobrym porządku element jest wyznaczony
przez swój odcinek poprzedników. Sprzeczność.

Rozważmy zbiór wszystkich izomorfizmów między odcinkami
początkowymi $U$ i $V$. Jest to \emph{zbiór}, ponieważ
składa się z podzbiorów $U\times V$.
Z poprzedniego akapitu izomorfizmy te są zgodne na
częściach wspólnych, więc ich suma $f$ jest izomorfizmem
między pewnymi odcinkami początkowymi $I\subseteq U$
i $J\subseteq V$. Jeśli oba byłyby właściwe,
wybierzmy najmniejsze elementy $u\in U\setminus I$
i $v\in V\setminus J$. Rozszerzenie $f\cup\{(u,v)\}$
byłoby jeszcze izomorfizmem odcinków początkowych,
wbrew definicji $f$ jako sumy wszystkich. Stąd
$I=U$ lub $J=V$. Jeśli oba są całymi zbiorami,
mamy izomorfizm $U\cong V$. W pozostałych przypadkach
jeden zbiór jest izomorficzny z właściwym odcinkiem
drugiego. Te trzy możliwości są rozłączne.
Rzeczywiście, żaden zbiór dobrze uporządkowany $W$
nie jest izomorficzny z własnym właściwym odcinkiem
początkowym $W_{<w}$.
Dla dowolnego ściśle rosnącego zanurzenia porządkowego
$f:W\to W$ pokażemy przez indukcję po dobrym porządku,
że $x\le f(x)$ dla każdego $x\in W$.
Wystarczy zauważyć, że dla wszystkich $y<x$
hipoteza indukcyjna daje $y\le f(y)<f(x)$,
zatem $f(x)$ jest większy od wszystkich $y<x$.
Gdyby $f(x)<x$, sam element $f(x)$
należałby do $W_{<x}$, a przy $y=f(x)$ hipoteza indukcyjna
wymusza $f(y)\ge y=f(x)$.
Z drugiej strony $y<x$ i ścisła monotoniczność
$f$ dają $f(y)<f(x)$, sprzeczność. Stąd $x\le f(x)$.
Gdyby obraz $f[W]$ zawierał się w $W_{<w}$,
mielibyśmy $w\le f(w)<w$, sprzeczność.
Nie wykorzystaliśmy aksjomatu wyboru.
\end{proof}

\section{Twierdzenie Hartogsa bez używania liczb porządkowych}
Do wykazania, że AC daje dobre uporządkowanie \emph{każdego}
zbioru, potrzebujemy specjalnego zbioru dobrze uporządkowanego,
którego nie da się różnowartościowo odwzorować do wyjściowego
zbioru. Zwykle opisuje się go za pomocą liczb porządkowych.
Ponieważ zdefiniujemy je później, użyjemy tu klas
izomorfizmu dobrych porządków; odpowiednia teoria ilorazów
została wprowadzona w rozdziale~\ref{ch:equivalence}.

\begin{twierdzenie}[Twierdzenie Hartogsa]\label{thm:choice-hartogs}
Dla każdego zbioru $A$ istnieje zbiór dobrze uporządkowany
$(H(A),<)$, dla którego nie istnieje iniekcja $H(A)\to A$.
Dowód nie wymaga aksjomatu wyboru.
\end{twierdzenie}
\begin{proof}
Niech $\mathcal W$ będzie zbiorem wszystkich par $(S,R)$,
gdzie $S\subseteq A$, a $R\subseteq S\times S$ jest dobrym
porządkiem na $S$. Ponieważ $R\subseteq S\times S\subseteq A\times A$,
wszystkie rozważane pary należą do
$\mathcal P(A)\times\mathcal P(A\times A)$.
Zbiór $\mathcal W$ jest jego podzbiorem określonym
przez formułę wyrażającą, że $R$ jest dobrym porządkiem na $S$;
istnieje więc przez separację
(Definicje~\ref{def:zf-power} i~\ref{def:zf-separation}
oraz Twierdzenie~\ref{thm:zf-product}).
Definiujemy relację równoważności: dwa takie porządki
są równoważne, jeśli istnieje między nimi izomorfizm
porządkowy. Jest to formuła ZF z parametrami:
izomorfizmy są podzbiorami produktów zbiorów uporządkowanych.
Dla każdego $w\in\mathcal W$ jego klasa $[w]$
istnieje przez separację w $\mathcal W$.
Zbiór wszystkich klas $H(A)$ istnieje przez zastępowanie
zastosowane do jednoznacznego przyporządkowania
$w\mapsto[w]$; można go też otrzymać przez separację
w $\mathcal P(\mathcal W)$.

Oznaczając przez $U,V$ reprezentantów z $\mathcal W$
wraz z ich relacjami dobrego porządku, dla dwóch klas
$[U],[V]\in H(A)$ piszemy $[U]<[V]$,
jeżeli $U$ jest izomorficzny z \emph{właściwym} odcinkiem
początkowym $V$. Jest to określenie niezależne
od wyboru przedstawicieli: izomorfizm porządków przenosi
właściwe odcinki początkowe na właściwe odcinki.
Lemat~\ref{lem:choice-wellorder-comparison} pokazuje,
że jest to ścisły porządek liniowy. Jest on dobry:
dla niepustego $K\subseteq H(A)$ weźmy $h\in K$.
Jeżeli nie ma elementu $K$ mniejszego od $h$, to $h$
jest najmniejszy. W przeciwnym razie podzbiór
$K\cap H(A)_{<h}$ ma element najmniejszy, ponieważ
odcinek $H(A)_{<h}$ jest izomorficzny z dobrze
uporządkowanym przedstawicielem klasy $h$:
elementowi $u$ przedstawiciela przypisujemy klasę
porządkową jego odcinka poprzedników. To przyporządkowanie
jest iniekcją: dwa różne odcinki początkowe jednego
dobrego porządku nie są izomorficzne. Jest surjekcją
na $H(A)_{<h}$ z definicji relacji $<$ i porównania
dobrych porządków. Zachowuje również porządek, gdyż
odcinek poprzedników wcześniejszego elementu jest
właściwym odcinkiem odcinka poprzedników późniejszego.

Przypuśćmy, że istnieje iniekcja $j:H(A)\to A$.
Przenieśmy porządek $H(A)$ na podzbiór $j[H(A)]\subseteq A$.
Ten podzbiór, z przeniesionym porządkiem, ma klasę
$h\in H(A)$, zatem całe $H(A)$ jest izomorficzne
z odcinkiem $H(A)_{<h}$. To niemożliwe dla dobrego
porządku, jak wykazaliśmy w Lemacie~\ref{lem:choice-wellorder-comparison}.
Wniosek: taka iniekcja nie istnieje.
\end{proof}

\begin{przyklad}[Zbiór Hartogsa dla dwóch elementów]
Jeśli $A=\{a,b\}$, to podzbiory $A$ mają tylko
moce $0$, $1$ albo $2$. Każdy dobrze uporządkowany
podzbiór ma więc jeden z trzech typów porządkowych:
pusty, jednoelementowy albo dwuelementowy.
Zatem $H(A)$ składa się z \emph{trzech} elementów,
uporządkowanych w kolejności odpowiadającej długościom
$0<1<2$. Nie istnieje iniekcja trzyelementowego
$H(A)$ do dwuelementowego $A$. Twierdzenie Hartogsa
rozszerza ten pomysł również na zbiory nieskończone,
dla których nie wystarcza liczenie elementów.
\end{przyklad}

\section{AC implikuje dobre uporządkowanie}
\begin{twierdzenie}[Twierdzenie o dobrym uporządkowaniu Zermela]
\label{thm:choice-wellorder}
W ZF aksjomat wyboru implikuje zasadę dobrego uporządkowania.
\end{twierdzenie}
\begin{proof}
Niech $A$ będzie dowolnym zbiorem. Dla $A=\varnothing$
teza jest oczywista; załóżmy $A\ne\varnothing$.
Stosujemy AC do rodziny wszystkich niepustych podzbiorów $A$:
istnieje $c:\mathcal P(A)\setminus\{\varnothing\}\to A$
taka, że $c(B)\in B$.
Niech $H:=H(A)$ z Twierdzenia~\ref{thm:choice-hartogs}.

Zbierzmy wszystkie funkcje $f:I\to A$, gdzie $I$ jest
odcinkiem początkowym $H$, i które spełniają dla każdego
$t\in I$ warunek
\[
 f(t)=c\bigl(A\setminus f[I_{<t}]\bigr).
\tag{$*$}
\]
Wymagamy przy tym, aby argument funkcji $c$ był niepusty.
Wszystkie takie funkcje tworzą zbiór: jest to
podzbiór $\mathcal P(H\times A)$ wyznaczony przez
formułę o wykresie funkcji, odcinku początkowym
i równaniu ($*$), z parametrami $H,A,c$.
Powstaje przez separację. Rodzina nie jest pusta,
ponieważ zawiera funkcję pustą o dziedzinie
pustego odcinka początkowego.
Dwie takie funkcje muszą zgadzać się na wspólnej dziedzinie:
gdyby nie, w najmniejszym miejscu rozbieżności miałyby
identyczne wcześniejsze wartości i przez ($*$) dostałyby
tę samą nową wartość. Ich suma $f$ jest więc funkcją
o dziedzinie $I$ będącej odcinkiem początkowym $H$
i nadal spełnia ($*$).

Gdyby $A\setminus f[I]$ było niepuste, a $I\ne H$,
można byłoby przedłużyć $f$ o najmniejszy element
$H\setminus I$, przypisując mu $c(A\setminus f[I])$.
To przeczyłoby temu, że $f$ jest sumą \emph{wszystkich}
dopuszczalnych funkcji. Natomiast jeśli $I=H$,
to $f$ jest iniekcją $H\to A$, gdyż każde nowe
wybrane $f(t)$ leży poza zbiorem wcześniejszych wartości,
co jest sprzeczne z twierdzeniem Hartogsa. Wobec tego
$A\setminus f[I]=\varnothing$, czyli $f$ jest bijekcją
$I\to A$. Przenosimy dobry porządek $I$ przez $f$
na $A$ i kończymy dowód.
\end{proof}

\section{Lemat Kuratowskiego--Zorna i jego dowód}
\begin{definicja}[Łańcuch i element maksymalny]\label{def:choice-chain}
W częściowo uporządkowanym zbiorze $(P,\le)$
\emph{łańcuchem} jest podzbiór $C\subseteq P$, w którym
dowolne dwa elementy są porównywalne. \emph{Ograniczenie
górne} łańcucha $C$ to element $u\in P$ spełniający
$\forall c\in C\ (c\le u)$. Element $m\in P$ jest
\emph{maksymalny}, jeśli nie istnieje $p\in P$ z $m<p$.
Są to pojęcia z rozdziału~\ref{ch:orders};
\emph{maksymalny} nie musi oznaczać \emph{największy}.
\end{definicja}

\begin{twierdzenie}[Lemat Kuratowskiego--Zorna]\label{thm:choice-zorn}
Niech $(P,\le)$ będzie \emph{zbiorem częściowo uporządkowanym}.
Załóżmy, że \emph{każdy} łańcuch $C\subseteq P$
ma ograniczenie górne \emph{należące do $P$}.
Założenie dotyczy również łańcucha pustego $C=\varnothing$.
Wówczas istnieje element maksymalny $m\in P$, czyli
\[
\exists m\in P\quad\forall p\in P\;(m\le p\ \Longrightarrow\ p=m).
\]
Nie żądamy istnienia elementu \emph{największego},
który musiałby być większy lub równy \emph{każdemu} $p\in P$.
\end{twierdzenie}
\begin{proof}
Założymy zasadę dobrego uporządkowania,
którą uzyskaliśmy z AC w Twierdzeniu~\ref{thm:choice-wellorder}.
Rozdzielmy argument na pięć kroków.

\emph{Krok 1: co właściwie chcemy skonstruować?}
Zdobędziemy łańcuch $C$ \emph{maksymalny względem zawierania}:
nie ma w $P$ żadnego większego łańcucha $C'\supsetneq C$.
Gdy go już znajdziemy, z założenia twierdzenia dostaniemy
górne ograniczenie $u\in P$ dla $C$.
Pokażemy wtedy, że $u$ jest szukanym elementem maksymalnym $P$.

\emph{Krok 2: kolejność rozpatrywania elementów.}
Z założenia lematu Zorna zbiór $P$ jest niepusty:
łańcuch pusty musi mieć ograniczenie górne w $P$.
Ustalmy dowolny dobry porządek $\triangleleft$ na $P$.
Nie musimy znać jego konkretnej postaci; wiemy jedynie,
że każdy niepusty podzbiór $P$ ma względem niego
pierwszy element. Porządek $\triangleleft$ służy tylko
do \emph{wyboru kolejności} i nie należy go mylić
z oryginalnym porządkiem częściowym $\le$.

\emph{Krok 3: przyjmowanie i odrzucanie kandydatów.}
Przechodzimy przez elementy $p\in P$ w kolejności
$\triangleleft$. Oznaczamy $p$ jako \emph{przyjęty},
jeśli dla każdego wcześniej przyjętego $q$
zachodzi $p\le q$ lub $q\le p$.
W przeciwnym razie odrzucamy $p$.
Na pierwszym etapie nie ma wcześniejszych kandydatów,
więc pierwszy element zostaje przyjęty.

Dla ścisłości: rozpatrujemy funkcje
$s:I\to\{0,1\}$, których dziedziny $I$ są odcinkami
początkowymi porządku $\triangleleft$, i które spełniają
\[
s(p)=1\ \Longleftrightarrow\
\forall q\in I\,(q\triangleleft p\land s(q)=1
\ \Longrightarrow\ (q\le p\lor p\le q)).
\]
Rodzina wszystkich takich funkcji jest zbiorem:
otrzymujemy ją przez separację
w $\mathcal P(P\times\{0,1\})$.
Zawiera funkcję pustą, która spełnia regułę
na pustym odcinku początkowym.
Dwie funkcje spełniające tę regułę muszą zgadzać się
na wspólnej dziedzinie: w pierwszym punkcie potencjalnej
różnicy mają identyczne wcześniej przyjęte elementy,
a zatem reguła narzuca tę samą wartość $0$ lub $1$.
Suma wszystkich takich funkcji jest zatem funkcją;
jej dziedzina jest odcinkiem początkowym.
Gdyby nie obejmowała całego $P$, moglibyśmy dodać
najmniejszy pominięty element (w porządku $\triangleleft$),
wyznaczając mu wartość przez powyższą regułę.
Przeczyłoby to temu, że wzięliśmy sumę wszystkich
możliwych poprawnych funkcji.
Otrzymujemy więc funkcję $s:P\to\{0,1\}$.

\emph{Krok 4: przyjęte elementy tworzą maksymalny łańcuch.}
Niech $C=\{p\in P:s(p)=1\}$.
Jeśli $p,q\in C$ i $p\triangleleft q$, to w chwili
przyjmowania $q$ element $p$ był już przyjęty.
Stąd $p\le q$ lub $q\le p$. Zatem $C$ jest łańcuchem.

Załóżmy, że istnieje większy łańcuch $C'$ z
$C\subsetneq C'$. Wybierzmy $v\in C'\setminus C$.
Każdy element $q\in C$ jest porównywalny z $v$,
bo $v,q\in C'$ i $C'$ jest łańcuchem.
Dotyczy to w szczególności każdego elementu
przyjętego \emph{przed} rozpatrzeniem $v$.
Reguła nakazywałaby zatem przyjąć $v$,
czyli $v\in C$. Sprzeczność. Łańcuch $C$
jest maksymalny względem zawierania.

\emph{Krok 5: od maksymalnego łańcucha do elementu maksymalnego.}
Z założenia twierdzenia istnieje $u\in P$ takie, że
$c\le u$ dla każdego $c\in C$.
Wtedy $C\cup\{u\}$ też jest łańcuchem:
element $u$ jest porównywalny ze wszystkimi $c\in C$.
Maksymalność $C$ wymusza $u\in C$.

Przypuśćmy teraz, że $u$ nie jest elementem maksymalnym
w $P$. Istniałby element $v\in P$ z $u<v$.
Ponieważ $c\le u<v$ dla każdego $c\in C$,
zbiór $C\cup\{v\}$ byłby łańcuchem ściśle
większym od $C$. Sprzeczność.
Zatem $u$ jest maksymalny w $P$.
\end{proof}

\begin{przyklad}[Dwie moce można porównać za pomocą iniekcji]
Niech $A,B$ będą zbiorami. Niech $P$ będzie zbiorem
wszystkich \emph{częściowych iniekcji} $f:D\to B$, gdzie
$D\subseteq A$, a porządkiem będzie przedłużanie funkcji.
Zbiór $P$ istnieje przez separację w $\mathcal P(A\times B)$:
każdy jego element jest wykresem funkcji różnowartościowej
o dziedzinie zawartej w $A$.
Wyjaśnijmy najpierw, dlaczego każdy łańcuch $C\subseteq P$
ma ograniczenie górne: suma $u=\bigcup C$ jest funkcją,
bo elementy łańcucha zgadzają się na wspólnych dziedzinach.
Jest też iniekcją: jeżeli $u(a)=u(a')$, to obie wartości
są zadane przez pewne funkcje z $C$; jedna z nich
przedłuża drugą, więc obie pary należą do tej samej
funkcji różnowartościowej. Stąd $a=a'$.
Wynika więc $u\in P$, a przy $C=\varnothing$ górnym
ograniczeniem jest funkcja pusta.

Lemat Zorna daje maksymalną częściową iniekcję
$f:D\to B$. Gdyby równocześnie $A\setminus D\ne\varnothing$
i $B\setminus f[D]\ne\varnothing$, moglibyśmy wybrać
po jednym elemencie $a\in A\setminus D$ i
$b\in B\setminus f[D]$ i przedłużyć $f$ przez
$f(a):=b$. Nadal mielibyśmy iniekcję, wbrew
maksymalności. Zatem albo $D=A$, i mamy iniekcję
$A\to B$, albo $f[D]=B$, i funkcja odwrotna $f^{-1}$
jest iniekcją $B\to D\subseteq A$.

Wniosek jest nieintuicyjnie silny: przy założeniu AC
dowolne dwa zbiory można porównać pod względem
istnienia iniekcji w jedną albo w drugą stronę.
Bez AC taka porównywalność nie musi zachodzić.
\end{przyklad}


\begin{przyklad}[Co naprawdę daje lemat Zorna?]
Weźmy $P=\{a,b\}$, w którym jedynymi porównaniami
są $a\le a$ i $b\le b$. Elementy $a,b$ są
nieporównywalne. Wszystkie łańcuchy to
$\varnothing$, $\{a\}$ i $\{b\}$.
Łańcuch pusty ma ograniczenie górne, np. $a$;
łańcuchy jednoelementowe mają ograniczenia
odpowiednio $a$ i $b$. Zatem założenia Zorna są spełnione.
Oba elementy $a,b$ są \emph{maksymalne},
ale żaden nie jest \emph{największy}.
Gdyby ktoś wywnioskował z lematu Zorna istnienie
elementu największego, ten przykład natychmiast
obaliłby taki wniosek.
\end{przyklad}

\section{Zasada maksymalności Hausdorffa}\label{sec:choice-hausdorff}

W dowodzie Twierdzenia~\ref{thm:choice-zorn}
zaczęliśmy od konstruowania maksymalnego \emph{łańcucha},
a dopiero potem znaleźliśmy element maksymalny.
Ten pośredni krok jest treścią osobnej zasady,
noszącej nazwisko Hausdorffa. Najpierw dokładnie
rozróżnijmy, co znaczy maksymalność \emph{łańcucha}
i maksymalność \emph{elementu} porządku.

\begin{definicja}[Łańcuch maksymalny]\label{def:hausdorff-maximal-chain}
Niech $(P,\le)$ będzie zbiorem częściowo uporządkowanym.
Łańcuch $C\subseteq P$ nazywamy \emph{maksymalnym}, jeśli
\[
\forall D\subseteq P\quad
 (D\text{ jest łańcuchem i }C\subseteq D\ \Longrightarrow\ D=C).
\]
Jest to maksymalność w porządku \emph{zawierania łańcuchów},
a nie istnienie jednego największego elementu w $C$.
\end{definicja}

\begin{definicja}[Zasada maksymalności Hausdorffa]\label{def:hausdorff-principle}
\emph{Zasada maksymalności Hausdorffa} głosi:
\emph{każdy zbiór częściowo uporządkowany ma
łańcuch maksymalny względem zawierania}.
Dla $P=\varnothing$ jedynym łańcuchem jest
$\varnothing$ i jest on maksymalny. Nie zakładamy
istnienia ograniczeń górnych łańcuchów.
\end{definicja}

\begin{przyklad}[Dwa różne łańcuchy maksymalne]
W zbiorze $P=\{a,b,c\}$ przyjmijmy $a<c$ i $b<c$,
a $a,b$ niech będą nieporównywalne.
Wówczas $C_1=\{a,c\}$ oraz $C_2=\{b,c\}$
są łańcuchami maksymalnymi. Nie możemy do $C_1$
dołożyć $b$, ponieważ $a$ i $b$ nie dają się porównać.
Analogicznie nie możemy do $C_2$ dołożyć $a$.
Łańcuch $\{c\}$ \emph{nie} jest maksymalny:
możemy powiększyć go do $C_1$.
Przykład pokazuje, że łańcuch maksymalny
nie musi być jednoznaczny.
\end{przyklad}

\begin{twierdzenie}[Równoważność Hausdorffa i Zorna nad ZF]\label{thm:hausdorff-zorn-equivalent}
W teorii ZF następujące dwa zdania są równoważne:
\begin{enumerate}[label=(\roman*)]
\item zasada maksymalności Hausdorffa;
\item lemat Kuratowskiego--Zorna.
\end{enumerate}
\end{twierdzenie}
\begin{proof}
\emph{Kierunek Zorn $\Rightarrow$ Hausdorff.}
Niech $(P,\le)$ będzie \emph{dowolnym} zbiorem
częściowo uporządkowanym. Zdefiniujmy
\[
\mathcal C:=\{C\subseteq P:C\text{ jest łańcuchem}\}.
\]
Rodzina $\mathcal C$ jest zbiorem, ponieważ jest podzbiorem
$\mathcal P(P)$ określonym przez własność łańcucha.
Jest także niepusta: $\varnothing\in\mathcal C$.
Uporządkujmy ją przez zawieranie $\subseteq$.
Chcemy zastosować Zorna do $(\mathcal C,\subseteq)$,
więc musimy wykazać jego założenie \emph{dla każdego
łańcucha w $\mathcal C$}.

Niech $\mathcal D\subseteq\mathcal C$ będzie łańcuchem
względem zawierania. To znaczy, że dla dowolnych
$C_1,C_2\in\mathcal D$ zachodzi
$C_1\subseteq C_2$ lub $C_2\subseteq C_1$.
Połóżmy $U:=\bigcup\mathcal D$.
Dla dowolnych $x,y\in U$ istnieją łańcuchy
$C_x,C_y\in\mathcal D$ takie, że $x\in C_x$ i $y\in C_y$.
Ponieważ $\mathcal D$ jest łańcuchem w porządku zawierania,
jeden z tych dwóch łańcuchów zawiera drugi.
Zatem zarówno $x$, jak i $y$ należą do jednego
łańcucha z $\mathcal D$, więc są porównywalne.
Wynika stąd $U\in\mathcal C$.
Ponadto $C\subseteq U$ dla każdego $C\in\mathcal D$,
czyli $U$ jest \emph{ograniczeniem górnym} $\mathcal D$
w zbiorze $\mathcal C$. Dla $\mathcal D=\varnothing$
mamy $U=\varnothing\in\mathcal C$, więc warunek
obejmuje również łańcuch pusty.

Stosujemy Zorna i otrzymujemy element
maksymalny $C_\ast\in\mathcal C$ względem zawierania.
Z definicji jest to dokładnie łańcuch maksymalny
w $P$. Otrzymaliśmy zasadę Hausdorffa.

\emph{Kierunek Hausdorff $\Rightarrow$ Zorn.}
Niech $(P,\le)$ spełnia założenie lematu Zorna:
każdy łańcuch ma ograniczenie górne należące do $P$.
Z zasady Hausdorffa otrzymujemy łańcuch maksymalny
$C\subseteq P$. Według założenia Zorna istnieje
$u\in P$ takie, że $c\le u$ dla każdego $c\in C$.
Rozważmy $C\cup\{u\}$.
Wszystkie pary elementów $C$ są porównywalne,
a każdy $c\in C$ spełnia $c\le u$.
Dlatego $C\cup\{u\}$ jest łańcuchem.
Z maksymalności $C$ wynika $u\in C$.

Gdyby $u$ nie był maksymalny w $P$, znaleźlibyśmy
$v\in P$ z $u<v$. Z przechodniości wynika
$c\le u<v$ dla każdego $c\in C$, a zatem
$C\cup\{v\}$ byłby łańcuchem właściwie
zawierającym $C$. To sprzeczne z maksymalnością $C$.
Wnioskujemy, że $u$ jest maksymalnym elementem $P$.
Dowiedliśmy lematu Zorna.
\end{proof}

\begin{uwaga}[Co odróżnia zasady Hausdorffa i Zorna?]
Zasada Hausdorffa mówi o \emph{maksymalnym łańcuchu}
bez zakładania, że łańcuchy mają ograniczenia górne.
Lemat Zorna mówi o \emph{maksymalnym elemencie}
pod dodatkowym założeniem istnienia takich ograniczeń.
To nie są identyczne zdania, ale właśnie udowodniliśmy,
że nad ZF mają dokładnie tę samą siłę.
\end{uwaga}

\section{Powrót od lematu Zorna do aksjomatu wyboru}
\begin{twierdzenie}[Równoważność głównych zasad]\label{thm:choice-equivalence}
W teorii ZF następujące cztery zdania są równoważne:
(i) aksjomat wyboru z Definicji~\ref{def:choice};
(ii) każde zbiory można dobrze uporządkować;
(iii) lemat Kuratowskiego--Zorna;
(iv) zasada maksymalności Hausdorffa z Definicji~\ref{def:hausdorff-principle}.
\end{twierdzenie}
\begin{proof}
Dowody (i)$\Rightarrow$(ii) i (ii)$\Rightarrow$(iii)
przeprowadziliśmy odpowiednio w Twierdzeniach
\ref{thm:choice-wellorder} i~\ref{thm:choice-zorn}.
Wystarczy wykazać (iii)$\Rightarrow$(i);
równoważność (iii)$\Leftrightarrow$(iv) została udowodniona
w Twierdzeniu~\ref{thm:hausdorff-zorn-equivalent}.

Niech $(A_i)_{i\in I}$ będzie dowolną rodziną niepustych
zbiorów. Rozważmy zbiór $P$ wszystkich \emph{częściowych
funkcji wyboru}: do $P$ należy funkcja $f:D\to\bigcup_{i\in I}A_i$
z dziedziną $D\subseteq I$, jeżeli $f(i)\in A_i$
dla każdego $i\in D$. Jest to zbiór istniejący przez
separację w $\mathcal P(I\times\bigcup_{i\in I}A_i)$:
warunek bycia funkcją i warunek $f(i)\in A_i$
są wyrażalne w języku ZF.
Uporządkujmy $P$ przez
przedłużanie: $f\le g$ wtedy i tylko wtedy, gdy
dziedzina $g$ zawiera dziedzinę $f$ i $g$ zgadza się
z $f$ na jej dziedzinie. Funkcja pusta należy do $P$.

Jeśli $C\subseteq P$ jest łańcuchem, to $u:=\bigcup C$
jest funkcją: dwie funkcje porównywalne przez
przedłużanie nie przypiszą różnych wartości temu samemu
indeksowi. Ponadto $u(i)\in A_i$ dla każdego
$i\in\dom u$, zatem $u\in P$ i $u$ jest
ograniczeniem górnym $C$ (także dla $C=\varnothing$).
Lemat Zorna daje maksymalną funkcję $f\in P$.
Gdyby $\dom f\ne I$, można byłoby wybrać \emph{jeden}
indeks $i_0\in I\setminus\dom f$ i \emph{jeden}
element $a_0\in A_{i_0}$, a następnie przedłużyć
$f$ o parę $(i_0,a_0)$. To przeczyłoby maksymalności.
Zatem $\dom f=I$ i $f$ jest funkcją wyboru.
\end{proof}

\begin{uwaga}[Dlaczego nie jest to dowód wyboru bez wyboru?]
Dowodząc AC z Zorna, \emph{zakładaliśmy} lemat Zorna.
Dowodząc Zorna z zasady dobrego uporządkowania,
\emph{zakładaliśmy} istnienie dobrego porządku na całym $P$.
Dowodząc zasady dobrego uporządkowania, zaczęliśmy
od AC. Nie wyprowadziliśmy więc żadnej z tych zasad
z samej ZF; wykazaliśmy jedynie, że w obecności ZF
każdą z nich można zastąpić dwiema pozostałymi.
\end{uwaga}

\begin{przyklad}[Równoległa i warunkowa postać wyboru]
Dla rodziny $(A_i)_{i\in I}$ \emph{iloczynem kartezjańskim rodziny}
nazywamy zbiór
\[
\prod_{i\in I}A_i:=\{c:I\to\bigcup_{i\in I}A_i:
\forall i\in I\ (c(i)\in A_i)\}.
\]
Zbiór tych funkcji istnieje przez separację
w $\mathcal P(I\times\bigcup_{i\in I}A_i)$;
dla $I=\varnothing$ zawiera funkcję pustą.
Czasem aksjomat wyboru przedstawia się w postaci:
„iloczyn kartezjański dowolnej rodziny niepustych
zbiorów jest niepusty”. Rzeczywiście, jego element
jest dokładnie funkcją wyboru $c$.
Ta równoważność wynika z samej definicji iloczynu
rodziny; nie należy mylić jej z iloczynem
\emph{dwóch} zbiorów, który istnieje w ZF bez AC.
\end{przyklad}

\section{Kiedy zasada wyboru jest rzeczywiście potrzebna?}

\begin{twierdzenie}[Wybór ze skończonej rodziny w ZF]
\label{thm:choice-finite-family}
Niech $I$ będzie zbiorem skończonym,
a $(A_i)_{i\in I}$ rodziną zbiorów niepustych.
W ZF istnieje funkcja $c$ o dziedzinie $I$
spełniająca $c(i)\in A_i$ dla wszystkich $i\in I$.
Teza obejmuje także $I=\varnothing$.
\end{twierdzenie}
\begin{proof}
Najpierw udowodnimy przez indukcję po $n\in\omega$
zdanie: \emph{każda rodzina $(B_k)_{k\in n}$
zbiorów niepustych ma funkcję wyboru
$d:n\to\bigcup_{k\in n}B_k$}.
Dla $n=0$ wystarcza funkcja pusta.
Załóżmy tezę dla $n$ i weźmy dowolną rodzinę
$(B_k)_{k\in S(n)}$ niepustych zbiorów.
Jej ograniczenie $(B_k)_{k\in n}$ ma
funkcję wyboru $d_n$ na mocy hipotezy indukcyjnej.
Ponieważ $B_n$ jest niepusty,
istnieje \emph{jeden} element $a\in B_n$.
Dołączamy parę $(n,a)$ do $d_n$.
Jest to funkcja na $S(n)=n\cup\{n\}$,
gdyż $n\notin n$
(Lemat~\ref{lem:zf-vn-successor}),
a każda jej wartość należy do odpowiedniego $B_k$.
Otrzymaliśmy tezę dla $S(n)$.

Teraz wróćmy do danej rodziny $(A_i)_{i\in I}$.
Ponieważ $I$ jest skończony, istnieją
$n\in\omega$ i bijekcja $e:n\to I$.
Zastosujmy przed chwilą dowiedzioną tezę
do rodziny $(B_k)_{k\in n}$ określonej przez
$B_k:=A_{e(k)}$. Otrzymujemy funkcję wyboru
$d:n\to\bigcup_{i\in I}A_i$.
Definiujemy $c(i):=d(e^{-1}(i))$ dla każdego
$i\in I$. Wtedy $c(i)\in A_i$.
W żadnym kroku nie zastosowaliśmy AC
do rodziny nieskończonej.
\end{proof}

\begin{przyklad}[Rodzina pusta i powtarzające się zbiory]
\label{ex:choice-empty-family}
Dla $I=\varnothing$ funkcja pusta
$\varnothing\to\varnothing$ jest funkcją wyboru.
Jeżeli $I=\{0,1,2\}$,
$A_0=A_1=\{a,b\}$ i $A_2=\{b\}$,
możemy wybrać $c(0)=a$, $c(1)=b$,
$c(2)=b$. Wartości funkcji wyboru
nie muszą być parami różne.
\end{przyklad}

\begin{twierdzenie}[Aksjomat wyboru a przekroje surjekcji]
\label{thm:choice-sections}
Nad ZF następujące warunki są równoważne:
\begin{enumerate}[label=(\roman*)]
\item obowiązuje aksjomat wyboru;
\item dla każdej surjekcji $p:E\to I$
istnieje funkcja $s:I\to E$ spełniająca
$p\circ s=\operatorname{id}_I$.
\end{enumerate}
Oba warunki obejmują puste zbiory.
\end{twierdzenie}
\begin{proof}
Załóżmy AC i niech $p:E\to I$ będzie surjekcją.
Dla $i\in I$ włókno
$E_i:=\{e\in E:p(e)=i\}$ jest niepuste.
Każde włókno istnieje przez separację,
a funkcja $i\mapsto E_i$ jest zbiorem przez
zastępowanie. Z AC istnieje
$s:I\to E$ takie, że $s(i)\in E_i$.
Zatem $p(s(i))=i$ dla każdego $i\in I$.

Odwrotnie, ustalmy rodzinę $(A_i)_{i\in I}$
zbiorów niepustych i jej sumę rozłączną
\[
E:=\{(i,a):i\in I,\ a\in A_i\}
\subseteq I\times\bigcup_{i\in I}A_i.
\]
Jest to zbiór zbudowany w ZF.
Rzut $p:E\to I$, $p(i,a)=i$,
jest surjekcją, ponieważ każde $A_i$ jest niepuste.
Przekrój $s$ tego rzutu musi mieć postać
$s(i)=(i,c(i))$ dla jednoznacznego
$c(i)\in A_i$. Schemat zastępowania daje
wykres funkcji $c$, a więc funkcję wyboru.
Dla $I=\varnothing$ zarówno suma rozłączna,
jak i przekrój są funkcjami pustymi.
\end{proof}

\begin{wniosek}[Transwersale relacji równoważności]
\label{cor:choice-transversals}
Nad ZF aksjomat wyboru jest równoważny temu,
że dla każdej relacji równoważności $\sim$
na zbiorze $X$ istnieje podzbiór $T\subseteq X$
zawierający dokładnie jeden element każdej klasy
równoważności ze zbioru ilorazowego $X/{\sim}$.
\end{wniosek}
\begin{proof}
Przy AC stosujemy Twierdzenie~\ref{thm:choice-sections}
do surjekcji ilorazowej $q:X\to X/{\sim}$.
Jej przekrój $s$ daje
$T:=s[X/{\sim}]$, przy czym
$q(s(C))=C$ dla każdej klasy $C$.
Zatem $T$ ma dokładnie jednego reprezentanta
z każdej klasy.

W drugą stronę dla dowolnej rodziny
$(A_i)_{i\in I}$ zbiorów niepustych
bierzemy sumę rozłączną
$E:=\{(i,a):i\in I,\ a\in A_i\}$
i relację $(i,a)\sim(j,b)$
wtedy i tylko wtedy, gdy $i=j$.
Jeżeli $T$ zawiera po jednym elemencie
każdej klasy, to jego jedyny element
we włóknie nad $i$ ma postać
$(i,c(i))$ z $c(i)\in A_i$.
Otrzymujemy funkcję wyboru $c$.
Dla $E=\varnothing$ transwersalą jest
zbiór pusty.
\end{proof}

\begin{przyklad}[Dlaczego suma musi być rozłączna?]
\label{ex:choice-disjoint-index}
Jeżeli $A_0=A_1=\{a\}$,
zwykła suma to tylko $\{a\}$,
która nie odnotowuje dwóch indeksów.
Suma rozłączna
$E=\{(0,a),(1,a)\}$ ma dwa elementy,
a jej włókna pamiętają, do którego
indeksu należy wybrany element.
\end{przyklad}

\section{Przykład zastosowania: istnienie bazy przestrzeni wektorowej}

\begin{twierdzenie}[Istnienie bazy z lematu Zorna]
\label{thm:choice-hamel-basis}
Niech $K$ będzie ciałem, a $V$ przestrzenią
wektorową nad $K$. Zakładając AC, istnieje
$B\subseteq V$ takie, że każdy $v\in V$
ma dokładnie jedno przedstawienie
\[
v=\sum_{b\in B}\lambda(b)b,
\]
gdzie $\lambda:B\to K$ jest funkcją
o skończonym nośniku
$\{b\in B:\lambda(b)\ne0\}$,
a sumowanie przebiega jedynie
po tym skończonym nośniku.
Zbiór $B$ nazywamy bazą przestrzeni $V$.
\end{twierdzenie}
\begin{proof}
Niech $P$ będzie zbiorem liniowo
niezależnych podzbiorów $C\subseteq V$.
Jest to zbiór otrzymany przez separację
w $\mathcal P(V)$, gdyż warunek
\[
\sum_{j=1}^r\lambda_jc_j=0
\Longrightarrow \lambda_1=\cdots=\lambda_r=0
\]
dla wszystkich skończonych list parami
różnych $c_j\in C$ i skalarów
$\lambda_j\in K$ można wyrazić w ZF.
Zbiór $P$ jest niepusty, bo $\varnothing\in P$.
Uporządkujmy go przez inkluzję.

Dla dowolnego łańcucha $\mathcal C\subseteq P$
zbiór $U:=\bigcup\mathcal C$ jest liniowo
niezależny. Istotnie, dowolne skończenie
wiele elementów $u_1,\ldots,u_r\in U$
należy łącznie do jednego zbioru
$C_*\in\mathcal C$: wybieramy
skończenie wiele zbiorów łańcucha
zawierających poszczególne $u_j$,
a spośród nich największy względem inkluzji.
Niezależność $C_*$ daje niezależność
tej skończonej rodziny. Dla pustego łańcucha
mamy $U=\varnothing\in P$.
Zatem każdy łańcuch ma górne ograniczenie w $P$.

Lemat Kuratowskiego--Zorna,
Twierdzenie~\ref{thm:choice-zorn},
daje maksymalny element $B\in P$.
Jeśli istniałby $v\in V$ poza
powłoką liniową $B$, to $B\cup\{v\}$
byłby liniowo niezależny.
Rzeczywiście, w każdej relacji
$\lambda v+\sum_{j=1}^r\mu_jb_j=0$
współczynnik $\lambda$ musi być zerowy,
gdyż w przeciwnym razie
$v=-\lambda^{-1}\sum_j\mu_jb_j$
należałby do powłoki liniowej $B$.
Przy $\lambda=0$ niezależność $B$
wymusza $\mu_j=0$ dla wszystkich $j$.
Byłoby to sprzeczne z maksymalnością $B$.
Stąd $B$ generuje $V$.

Jeżeli dwie skończone kombinacje
elementów $B$ przedstawiają ten sam wektor,
odejmujemy je i grupujemy współczynniki
przy tych samych wektorach.
Otrzymujemy skończoną relację liniową
równą zeru. Niezależność $B$
wymusza zerowanie wszystkich współczynników
różnicy, czyli jednoznaczność funkcji $\lambda$.
\end{proof}

\begin{przyklad}[Przestrzeń zerowa]\label{ex:choice-zero-space}
Jeśli $V=\{0\}$, zbiór pusty jest
liniowo niezależny i generuje $V$:
jego jedyna skończona kombinacja liniowa
to pusta suma, czyli wektor zerowy.
Zatem bazą $V$ jest $\varnothing$.
\end{przyklad}
